Sequential hypothesis testing offers flexibility over fixed-sample designs, but stopping
rules coupled to decision criteria risk confirmation bias through early peeking.
The HDI+ROPE algorithm exemplifies this: it stops as soon as the Bayesian posterior
Highest Density Interval (HDI) lies entirely inside or outside a pre-specified Region
of Practical Equivalence (ROPE), enabling positive \textit{acceptance} of the null, 
unlike Null Hypothesis Significance Testing. But without a precision requirement,
it can stop on imprecise posteriors, incurring a systematic false-positive cost.
Decoupled methods, such as ``Precision is the Goal'' (PitG), eliminate this bias by
halting only once a target HDI width $\omega_{\rm goal}$ is reached; yet PitG
frequently yields inconclusive outcomes rather than a premature verdict, most
critically when the null is true.
We propose ``Decisive Precision is the Goal'' (DPitG), which requires \textit{both}
the precision target and a conclusive verdict to be satisfied simultaneously before
stopping.
In fair coin experiments (${\rm ROPE} = \theta_{\rm null} \pm 0.05$), PitG's
conclusiveness ranges from $82.5\%$ to $0\%$ across $\omega_{\rm goal}=0.06$--$0.10$;
DPitG's remains essentially flat at $97.2\%$--$97.9\%$.
At $\omega_{\rm goal}=0.08$, DPitG reduces PitG's $63.3\%$ inconclusive rate to $2.3\%$
at a median cost of only $4.7\%$ more samples than PitG, while preserving
a zero false-positive rate; HDI+ROPE matches this conclusiveness only at a $6.2\%$
false-positive cost.
Results are demonstrated on single-group binary data; the framework extends naturally
to continuous outcomes and two-group comparisons.
This advantage holds broadly across true and null parameter values.
We provide a closed-form planning formula
($N_{\rm goal} \propto V/\omega_{\rm goal}^2$, where $V$ is the per-observation variance),
an online interactive calculator, and open-source code.
DPitG's fully pre-specified stopping rule is compatible with pre-registration standards
and is the method of choice whenever a reliable verdict is required and the budget
reaches $N_{\rm goal}$.

%%TODO: note that the phrasing of "nlike methods that can only reject the null hypothesis, the precision-based framework also permits its positive acceptance" are 
%% a bit decieinvg, b/c, e.g, Bayes Factors can also accept the null, but only when the data are extremely strong, whereas the precision-based framework can accept the null when the data are sufficiently precise, which is a more common occurrence.

%% -----------------------------------------------------------------------
%% IMRaD ALTERNATIVE (for reference)
%% -----------------------------------------------------------------------
% \textbf{Background.} Sequential hypothesis testing offers flexibility over
% fixed-sample designs, but risks confirmation bias when stopping rules are
% coupled to decision criteria.
%
% \textbf{Gap.} Existing decoupled approaches such as ``Precision is the Goal''
% (PitG) eliminate this bias but produce high rates of inconclusive outcomes,
% especially when the null hypothesis is true.
%
% \textbf{Method.} We propose ``Decisive Precision is the Goal'' (DPitG), which
% requires both a target posterior precision and a conclusive hypothesis outcome
% to be satisfied simultaneously before stopping.
%
% \textbf{Results.} In fair coin simulations ($\theta_{\rm true}=0.5$,
% $\omega_{\rm goal}=0.08$), PitG yields 63.3\% inconclusive outcomes; DPitG
% reduces this to 2.3\% with only a 4.7\% median increase in sample size,
% while preserving zero false positives.
%
% \textbf{Impact.} DPitG offers a practical and principled solution for sequential
% testing when decisive outcomes are required. Open-source code and an interactive
% calculator are provided.
%% -----------------------------------------------------------------------
